\documentclass[english, parskip=full]{scrartcl}
\usepackage[utf8]{inputenc}  % Kodierung der Datei
\usepackage[english]{babel}  % Sprache auf Deutsch 
\usepackage{color}
\usepackage{graphicx, gensymb}
\usepackage{amsfonts, cancel, amssymb}
\usepackage{amsmath, amsthm, array, esint}
\usepackage{booktabs, multirow}  % schönere Tabellen
\usepackage{caption}  % Anpassung der Captions
\usepackage{scrlayer-scrpage}    % Manipulation des Seitenstils
\pagestyle{scrheadings}  % Stil für die Seite setzen
\automark{section}  % Abschnittsnamen als Seitenbeschriftung 
\setheadsepline{.5pt}    % Kopzeile durch Linie abtrennen
\usepackage[hidelinks]{hyperref}  % Links und weitere PDF-Features
\usepackage{typearea, tikz, pgffor, verbatim}
\usetikzlibrary{calc, arrows.meta,arrows, graphs, quotes, positioning}
\usepackage[cache=false, outputdir=build]{minted}
\usemintedstyle{vs}
\usepackage{placeins} % provides \FloatBarrier

\definecolor{bg}{rgb}{0.9,0.9,0.9}
\newcommand\numberthis{\addtocounter{equation}{1}\tag{\theequation}}
\usepackage{svg}
\usepackage{siunitx}
\usepackage{physics}
\usepackage{tensor}
\usepackage[compat=1.1.0]{tikz-feynman}
\renewcommand{\bra}{}
\newcommand{\kom}{}
\newcommand{\brak}{}
\renewcommand{\braket}{}
\newcommand{\intx}{}
\newcommand{\intr}{}
\newcommand{\kla}{}
\newcommand{\klam}{}
\newcommand{\klammer}{}
\newcommand{\oper}{}
\newcommand{\tecxt}{}
\newcommand{\Bra}{}
\newcommand{\Ket}{}
\renewcommand{\i}{\text{i}}
\newcommand{\e}{\text{e}}
\newcommand*{\Fourier}{\textsc{Fourier}\ }
\newcommand*{\F}{\mathcal{F}}
\newcommand*{\sinc}{\text{sinc\,}}
\newcommand{\conv}{\mathop{\scalebox{3.5}{\raisebox{-0.38ex}{\makebox[0.5\width][c]{$\ast$}}}}\limits}

\newcommand*{\titel}{2 Self Energy and Coulomb interaction}
\newcommand*{\autor}{Joachim Schwardt}

\hypersetup{pdfauthor={\autor}, pdftitle={\titel}}

%\titlehead{titlehead}
\subject{Quantum Many-Body Physics}
\title{\titel}
\author{\autor}
\date{\today}

\begin{document}

%\begin{titlepage}
\maketitle  % Titel setzen
%\tableofcontents  % Inhaltsverzeichnis setzen
%\end{titlepage}

\section{Self-energy of the Fermi-Hubbard model}
%\feynmandiagram [large, vertical=e to f] {
%a -- [fermion] b -- [photon, momentum'=\(k\)] c -- [fermion] d,
%b -- [fermion, momentum=\(p_{1}\)] e -- [fermion, momentum=\(p_{2}\)] c,
%e -- [gluon]
%f,
%h -- [fermion] f -- [fermion] i;
%};
%\feynmandiagram [vertical=e to f]  {
%    a -- [fermion] b -- [photon, momentum=$k$] c -- [fermion] d,
%    b -- [fermion, momentum'=$p_1$] e -- [fermion, momentum'=$p_2$] c,
%    e -- [gluon] f,
%    g -- [fermion] f -- [fermion] i;
%};
Consider the Fermi-Hubbard model described by the Hamiltonian
\begin{align}
\hat{H} &= \sum_{jk\sigma} t_{jk} \hat{c}_{j\sigma}^\dagger \hat{c}_{k\sigma} + \frac{U}{2} \sum_{j\sigma} \hat{n}_{j\sigma} \hat{n}_{j,-\sigma}
\end{align}
where the $t_{jk}$ denote the (spin-independent) hopping amplitudes, $U > 0$ the Hubbard interaction strength and $\hat{n}_{j\sigma} = \hat{c}_{j\sigma}^\dagger \hat{c}_{j\sigma}$. 
Note that the Hubbard interaction only couples to opposite spins.
This is evident from the the interaction term containing $\hat{c}_{j\sigma} \hat{c}_{j,-\sigma}$, which requires the interacting states to have opposite spin.
Note that the Feynman rules may be expressed as
\begin{align}
\feynmandiagram [small, horizontal=i to f, baseline={-0.5ex}] {
    f [particle=$\sigma'$] -- [dashed, momentum'=$P$] i [particle=$\sigma$];
}; &= -\i U(P) \delta_{\sigma,-\sigma'} = -\i \delta_{\sigma, -\sigma'} \frac{U}{2},\\
\feynmandiagram [small, horizontal=a to b, layered layout] {
   a [nudge=(0:10mm)] -- b [dot] -- [out=135, in=45, loop, min distance=1cm] b,
 }; &= n_0(\mu),\\
\feynmandiagram [small, horizontal=i to f, baseline={-0.5ex}] {
    f [particle=$\sigma'$] -- [fermion, momentum'=$P$] i [particle=$\sigma$];
}; &= \i G_{0,\sigma\sigma'}^\tecxt\text{c}(P) = \i \delta_{\sigma\sigma'} G_0^\tecxt\text{c}(P) = \i\delta_{\sigma\sigma'} \frac{1}{\omega - \epsilon_\textbf{k}(1 - \i 0^+)},\\
\feynmandiagram [small, horizontal=i to f, baseline={-0.5ex}] {
    f [particle=$\sigma'$] -- [double, double distance=0.2ex, with arrow=0.5, arrow size=0.11em, momentum'=$P$] i [particle=$\sigma$];
}; &= \i G_{\sigma\sigma'}^\tecxt\text{c}(P).
\end{align}
The full propagator is
\begin{align}
\feynmandiagram [small, horizontal=i to f, baseline={-0.5ex}] {
    f [particle=$\sigma'$] -- [double, double distance=0.2ex, with arrow=0.5, arrow size=0.11em, momentum'=$P$] i [particle=$\sigma$];
}; &= \feynmandiagram [small, horizontal=i to f, baseline={-0.5ex}] {
    f [particle=$\sigma'$] -- [fermion, momentum'=$P$] i [particle=$\sigma$];
}; + \feynmandiagram [small, horizontal=i to f, baseline={-0.5ex}] {
    f [particle=$\sigma'$, nudge=(180:0.3cm)] -- [fermion, momentum'=$P$] a [dot, label=-90:$\sigma_1$] -- [fermion, momentum'=$P$] i [particle=$\sigma$, nudge=(0:0.3cm)],
     a -- [draw=none] c [nudge=(90:1.5cm)],
     a -- [scalar] b [dot, nudge=(-90:0.6cm), label=0:$\sigma_2$] -- [out=135, in=45, loop, min distance=1cm] b;
}; + \feynmandiagram [small, layered layout, horizontal=i to f, baseline={-0.5ex}] {
    f [particle=$\sigma'$] -- [fermion, momentum'=$P$] a [dot, label=-90:$\sigma_2$] -- [fermion, momentum=$P-Q$] b [dot, label=-90:$\sigma_1$] -- [fermion, momentum'=$P$] i [particle=$\sigma$],
    a -- [scalar, half right, momentum'=$Q$] b;
}; + \dots
\end{align}
The first diagram is of course just the free propagator.
The third diagram is also simple, because it contains a term $\sim \delta_{\sigma_1\sigma_2}$ from the free propagator but also a term $\sim \delta_{\sigma_1, -\sigma_2}$ from the Hubbard interaction and hence vanishes.
The second diagram gives
\begin{align}
\feynmandiagram [small, horizontal=i to f, baseline={-0.5ex}] {
    f [particle=$\sigma'$, nudge=(180:0.3cm)] -- [fermion, momentum'=$P$] a [dot, label=-90:$\sigma_1$] -- [fermion, momentum'=$P$] i [particle=$\sigma$, nudge=(0:0.3cm)],
     a -- [draw=none] c [nudge=(90:1.5cm)],
     a -- [scalar] b [dot, nudge=(-90:0.6cm), label=0:$\sigma_2$] -- [out=135, in=45, loop, min distance=1cm] b;
}; &= \sum_{\sigma_1\sigma_2} \i \delta_{\sigma\sigma_1} G_0^\tecxt\text{c}(P) \i \delta_{\sigma_1\sigma'} G_0^\tecxt\text{c}(P) n_0(\mu) \kla\left( -\i \frac{U}{2} \right) \delta_{\sigma_1, -\sigma_2} = \\
&= \i\delta_{\sigma\sigma'} \frac{U}{2} \big( G_0^\tecxt\text{c}(P) \big)^2 n_0(\mu).
\end{align}
The self-energy to first order in perturbation theory is given by
\begin{align}
-\i\Sigma &= \feynmandiagram [small, horizontal=i to f, baseline={-0.5ex}] {
    f [nudge=(180:1cm)] -- [draw=none] a [dot, label=-90:{$\sigma_1, \delta_{\sigma\sigma'}$}] -- [draw=none] i [nudge=(0:1cm)],
     a -- [draw=none] c [nudge=(90:1.5cm)],
     a -- [scalar] b [dot, nudge=(-90:0.6cm), label=0:$\sigma_2$] -- [out=135, in=45, loop, min distance=1cm] b;
}; + \feynmandiagram [small, layered layout, horizontal=i to f, baseline={-0.5ex}] {
    f -- [draw=none] a [dot, label=0:$\sigma_2$] -- [fermion, momentum=$P-Q$] b [dot, label=180:$\sigma_1$] -- [draw=none] i,
    a -- [scalar, half right, momentum'=$Q$] b;
}; + \mathcal{O}(U^2) \simeq \\
&\simeq \i\delta_{\sigma\sigma'} \frac{U}{2} n_0(\mu).
\end{align}
The Dyson equation is
\begin{align}
G^\tecxt\text{c} &= \kla\left( G_0^\tecxt\text{c}(P)^{-1} - \Sigma(P)^{-1} \right)^{-1}.
\end{align}
But since the spin structure is completely diagonal, we only need to consider the scalar part 
\begin{align}
G^\tecxt\text{c}(\omega,\textbf{k},\sigma) &= \kla\left( G_0^\tecxt\text{c}(\omega,\textbf{k},\sigma)^{-1} - \Sigma(\omega,\textbf{k},\sigma)^{-1} \right)^{-1} = \frac{1}{\omega - \epsilon_\textbf{k}(1 - \i 0^+) + \frac{U}{2} n_0(\mu)} + \mathcal{O}(U^2).
\end{align}
One of the second order diagrams for the self-energy that does not vanish due to the spin structure is
\begin{align*}
&\feynmandiagram [small, layered layout, horizontal=i to f, baseline={-0.5ex}] {
    f [particle=$\sigma'$] -- [fermion, momentum'=$P$] a[dot, label=90:$\sigma_1$] -- [fermion] b[dot, label=90:$\sigma_2$] -- [fermion, momentum'=$P$] i [particle=$\sigma$],
    a -- [scalar, quarter left, momentum=$Q_1$] c[dot, label=0:$\sigma_3$] -- [fermion, half right, momentum'=$Q_1 + Q_2$] d[dot, label=180:$\sigma_4$],
    d -- [fermion, half right, momentum'=$Q_2$] c,
    d -- [scalar, quarter left, momentum=$Q_1$] b;
}; = \\
&= (-1)^1 \sum_{\sigma_1 \sigma_2 \sigma_3 \sigma_4} \intr\int_{\mathbb{R}^{4}} \intr\int_{\mathbb{R}^{4}} \i\delta_{\sigma\sigma_2} G_0^\tecxt\text{c}(P) \i\delta_{\sigma\sigma_2} G_0^\tecxt\text{c}(P-Q_1) \i \delta_{\sigma_1\sigma'} G_0^\tecxt\text{c}(P) \times \\
&\hspace{4cm} \times \left( -\i \frac{U}{2} \right) \delta_{\sigma_1, -\sigma_3} \left( -\i \frac{U}{2} \right) \delta_{\sigma_2, -\sigma_4} \i\delta_{\sigma_3\sigma_4} G_0^\tecxt\text{c}(Q_1 + Q_2)  \times \\
&\hspace{4cm} \times \i\delta_{\sigma_3\sigma_4} G_0^\tecxt\text{c}(Q_2)  \, \frac{\mathrm{d}^{4}Q_1}{(2\pi)^4} \, \frac{\mathrm{d}^{4}Q_2}{(2\pi)^4}. \numberthis
\end{align*}
\section{Thomas-Fermi screening of the Coulomb interaction}
Consider the degenerate electron gas in three dimensions described by the Hamiltonian
\begin{align}
\hat{H} &= \sum_{\textbf{k}\sigma} \epsilon_\textbf{k} \hat{c}_{\textbf{k}\sigma}^\dagger \hat{c}_{\textbf{k}\sigma} + \frac{1}{2N} \sum_{\textbf{kk}'\sigma\sigma'} \sum_{\textbf{q}\neq 0} V_{\textbf{q}} \hat{c}_{\textbf{k} + \textbf{q} \sigma}^\dagger \hat{c}_{\textbf{k}' - \textbf{q}\sigma'}^\dagger \hat{c}_{\textbf{k}'\sigma'} \hat{c}_{\textbf{k}\sigma}
\end{align}
where $\epsilon_\textbf{k} = \frac{k^2}{2m} - \epsilon_F$, $V_\textbf{q} = \frac{1}{q^2}$ and $N$ denoting the total number of sites.
The vacuum polarization diagram in momentum-frequency space to lowest order in perturbation theory is given by 
\begin{align}
\i\Pi &= \feynmandiagram [layered layout, small, horizontal=a to b, baseline={-0.5ex}] {
    i [nudge=(0:5mm)] -- [small] a [dot] -- [fermion, half left, looseness=1.5, momentum={$k+q, \omega + \epsilon$}] b [dot] -- [fermion, half left, looseness=1.5, momentum={$q, \epsilon$}] a, 
    b -- f [nudge=(180:5mm)];
}; = \underbrace{(-1)^1}_{\tecxt\text{1 closed loop}} \sum_\sigma \intr\int_{\mathbb{R}^{4}} \i G_0^\tecxt\text{c}(\textbf{k} + \textbf{q}, \omega + \epsilon) \i G_0^\tecxt\text{c}(\textbf{q}, \epsilon) \, \frac{\mathrm{d}^{4}Q}{(2\pi)^4}.
\end{align}
Since $G_0^\tecxt\text{c}(\textbf{k},\omega) = \frac{1}{\omega - \epsilon_\textbf{k} + \tecxt\text{sign\,} \epsilon_\textbf{k} \i 0^+} = \frac{1}{\omega - \epsilon_\textbf{k} (1 - \i 0^+)}$ we find
\begin{align}
\Pi &= \underbrace{2}_\text{spin} \frac{1}{2\pi\i} \intr\int_{\mathbb{R}^{3}} \intr\int_{\mathbb{R}^{ }} \frac{1}{\omega + \epsilon - \epsilon_{\textbf{k} + \textbf{q}} (1 - \i 0^+)} \frac{1}{\epsilon - \epsilon_{\textbf{q}} (1 - \i 0^+)} \, \mathrm{d}^{ }\epsilon \, \frac{\mathrm{d}^{3}q}{(2\pi)^3}.
\end{align}
The integrand has poles at $\epsilon = \epsilon_{\textbf{q}}(1 - \i 0^+)$ and $\epsilon = \epsilon_{\textbf{k} + \textbf{q}} (1 - \i 0^+) - \omega$.
If both poles lie on the same half of the complex plane, we can close the contour on the opposing half and get a result of zero.
Thus the signs of $\epsilon_{\textbf{k} + \textbf{q}}$ and $\epsilon_\textbf{q}$ must be different, which is encapsulated by including the term $\Theta(\epsilon_{\textbf{k} + \textbf{q}}) \Theta(-\epsilon_\textbf{q})$.
Using the residue theorem and $\Theta(\epsilon) = n_0(\epsilon)$ for $T=0$ the integral becomes
\begin{align}
\Pi &= 2 \frac{1}{2\pi\i} \intr\int_{\mathbb{R}^{3}} 2\pi\i \klam\left[ \frac{\Theta(\epsilon_{\textbf{k} + \textbf{q}}) \Theta(-\epsilon_\textbf{q})}{\omega + \epsilon_\textbf{q} - \epsilon_{\textbf{k} + \textbf{q}} + \i 0^+} + \frac{\Theta( - \epsilon_{\textbf{k} + \textbf{q}}) \Theta( \epsilon_\textbf{q})}{-\omega + \epsilon_{\textbf{k} + \textbf{q}} - \epsilon_\textbf{q} + \i 0^+} \right]  \, \frac{\mathrm{d}^{3}q}{(2\pi)^3} = \\
&= 2 \intr\int_{\mathbb{R}^{3}} \left[ \frac{n_0(\epsilon_{\textbf{k} + \textbf{q}}) n_0( -\epsilon_\textbf{q})}{\omega + \epsilon_\textbf{q} - \epsilon_{\textbf{k} + \textbf{q}} + \i 0^+} - \frac{n_0(\epsilon_{\textbf{k} + \textbf{q}}) n_0( -\epsilon_\textbf{q})}{\omega + \epsilon_{\textbf{k} + \textbf{q}} - \epsilon_\textbf{q} - \i 0^+} \right] \, \frac{\mathrm{d}^{3}}{(2\pi)^3}.
\end{align}
Using the following identities 
\begin{align}
n_0(-\epsilon_\textbf{q}) &= \Theta(-\epsilon_\textbf{q}) = \Theta \kla\left( -\frac{q^2 - k_F^2}{2m} \right) = \Theta (k_F - q),\\
n_0(\epsilon_{\textbf{k} + \textbf{q}}) &= \Theta (\epsilon_{\textbf{k} + \textbf{q}}) = \Theta \kla\left( \frac{(\textbf{k} + \textbf{q})^2 - k_F^2}{2m} \right) = \Theta ( |\textbf{k} + \textbf{q}| - k_F)
\end{align}
we finally arrive at
\begin{align*}
\Pi(\textbf{k},\omega) &= 2 \intr\int_{\mathbb{R}^{3}} \Theta(|\textbf{k} + \textbf{q}| - k_F) \Theta(k_F - q) \times \\
&\hspace{2cm} \times \left[ \frac{1}{\omega + \epsilon_\textbf{q} - \epsilon_{\mathbf{k} + \textbf{q}} + \i 0^+} - \frac{1}{\omega - \epsilon_\textbf{q} + \epsilon_{\textbf{k} + \textbf{q}} - \i 0^+} \right] \,\frac{ \mathrm{d}^{3}q}{(2\pi)^3}. \numberthis
\end{align*}
%LINDHARD:
%\begin{align}
%\Pi(\textbf{k}, 0) &= -\frac{2mk_F}{(2\pi)^2} \klam\left[ 1 + \frac{k_F - (k/2)^2}{k_Fk} \log \left\vert \frac{k_F + k/2}{k_F - k/2} \right\vert \right].
%\end{align}
Using the substitution $k^2 + q^2 + 2kq\cos\theta = u$ with $\tecxt\text{d}u = 2kq \tecxt\text{d}(\cos\theta)$ we find
\begin{align}
\Pi(\textbf{k},0) &= \lim_{\alpha \rightarrow 0^+} \frac{2}{(2\pi)^2} \intx\int_{0}^{k_F} \intx\int_{-1}^{1} \frac{2\Theta\kla\left( \sqrt{k^2 + q^2 + 2kq\cos \theta} - k_F \right)}{\frac{q^2}{2m} - \frac{q^2 + k^2 + 2kq\cos\theta}{2m} + \i\alpha} \, \mathrm{d}(\cos\theta) \, q^2\mathrm{d}q = \\
&= \lim_{\alpha \rightarrow 0^+} \frac{8m}{(2\pi)^2} \intx\int_{0}^{k_F} \intx\int_{(k-q)^2}^{(k+q)^2} \frac{\Theta(\sqrt{u} - k_F)}{q^2 - u + \i\alpha} \, \frac{1}{2kq}\mathrm{d}u \, q^2\mathrm{d}q.
\end{align}
We first need to evaluate the Theta-function.
Since $0\le q,k\le k_F$ the lower bound satisfies $(k-q)^2 \le k_F^2$, so we may replace it by $k_F^2$.
However, this is only justified given $(k+q)^2 \ge k_F^2$.
This condition is equivalent to $q\ge k_F - k$, which allows us to write
%It is helpful to split the $q$-integration into two parts via
%\begin{align}
%\int_0^{k_F - k} \int_{(k-q)^2}^{(k+q)^2} &\Theta(\sqrt{u} - k_F) \dots + \int_{k_F - k}^{k_F} \intx\int_{(k-q)^2}^{(k+q)^2} \Theta(\sqrt{u} - k_F) \dots = \\
%&= 0 + \int_{k_F - k}^{k_F} \int_{k_F^2}^{(k+q)^2}\dots 
%\end{align}
%Integrating by parts yields
%\begin{align*}
%\Pi &= \lim_{\alpha \rightarrow 0^+} \frac{4m/k}{(2\pi)^2} \intx\int_{0}^{k_F} q \Theta(u-k_F) \log |q^2 - u + \i\alpha| \bigg\vert_{(k-q)^2}^{(k+q)^2} \mathrm{d}q\ - \\
%&\hspace{1.5cm} - \frac{4m/k}{(2\pi)^2} \int_0^{k_F} q\int_{(k-q)^2}^{(k+q)^2} \delta(u - k_F) \log |q^2 - u + \i\alpha|  \, \mathrm{d}u\, \mathrm{d}q 
%\end{align*}
%\begin{align}
%%%% BEGIN FOUR INTEGRALS
%\Pi(\textbf{k},0) &= \lim_{\alpha \rightarrow 0^+} \frac{4m/k}{(2\pi)^2} \intx\int_{k_F - k}^{k_F} \intx\int_{k_F^2}^{(k+q)^2} \frac{1}{q^2 - u + \i\alpha} \,\mathrm{d}u \, q\mathrm{d}q = \\
%&= \lim_{\alpha \rightarrow 0^+} \frac{4m/k}{(2\pi)^2} \intx\int_{k_F - k}^{k_F} \log \left\vert \frac{q^2 - (k+q)^2 + \i\alpha}{q^2 - k_F^2 + \i\alpha} \right\vert \, q\mathrm{d}q = \\
%&= \frac{4m/k}{(2\pi)^2} \intx\int_{k_F - k}^{k_F} \kla\left( \log |2k| + \log \left\vert q + \frac{k}{2} \right\vert  - \log |q - k_F| - \log |q + k_F|  \right) \, q\mathrm{d}q.
%\end{align}
%%%% END FOUR INTEGRALS
%%%% BEGIN TWO INTEGRALS
%\begin{align}
%\Pi(\textbf{k},0) &= \lim_{\alpha \rightarrow 0^+} \frac{4m/k}{(2\pi)^2} \intx\int_{k_F - k}^{k_F} \intx\int_{k_F^2}^{(k+q)^2} \frac{1}{q^2 - u + \i\alpha} \,\mathrm{d}u \, q\mathrm{d}q = \\
%&= \lim_{\alpha \rightarrow 0^+} \frac{4m/k}{(2\pi)^2} \intx\int_{k_F - k}^{k_F} \log \left\vert \frac{q^2 - (k+q)^2 + \i\alpha}{q^2 - k_F^2 + \i\alpha} \right\vert \, q\mathrm{d}q = \\
%&= \frac{4m/k}{(2\pi)^2} \intx\int_{k_F - k}^{k_F} \kla\left( \log |k(2q - k)| - \log |q^2 - k_F^2|  \right) \, q\mathrm{d}q = \\
%&= \frac{4m/k}{(2\pi)^2} \klam\left[ I_1 - I_2 \right].
%\end{align}
%%%% END TWO INTEGRALS
\begin{align}
\Pi(\textbf{k},0) &= \lim_{\alpha \rightarrow 0^+} \frac{4m/k}{(2\pi)^2} \intx\int_{k_F - k}^{k_F} \intx\int_{k_F^2}^{(k+q)^2} \frac{1}{q^2 - u + \i\alpha} \,\mathrm{d}u \, q\mathrm{d}q = \\
&= \lim_{\alpha \rightarrow 0^+} \frac{4m/k}{(2\pi)^2} \intx\int_{k_F - k}^{k_F} \log \left\vert \frac{q^2 - (k+q)^2 + \i\alpha}{q^2 - k_F^2 + \i\alpha} \right\vert \, q\mathrm{d}q = \\
&= \frac{4m/k}{(2\pi)^2} \intx\int_{k_F - k}^{k_F} q \log \left\vert \frac{k^2 + 2qk}{q^2 - k_F^2} \right\vert \,\mathrm{d}q.
\end{align}
%Note that 
%\begin{align}
%\int x\log x\, \mathrm{d}x &= \frac{x^2}{2} \log x - \int \frac{x^2}{2} \frac{1}{x}\, \mathrm{d}x = \frac{x^2}{2} \log x - \frac{x^2}{4},\\
%\intx\int_{ }^{ } x \log |x+b|\, \mathrm{d}x &= \intx\int_{ }^{ } (u-b) \log |u|\, \mathrm{d}u = \frac{u^2}{2} \log |u| - \frac{u^2}{4} - b u\log |u| + bu = \\
%&= \frac{u}{2} \klam\left[ (u - 2b) \log |u| - \frac{u}{2} + 2b \right] = \\
%&= \frac{x+b}{2} \klam\left[ (x-b) \log |x+b| - \frac{x+b}{2} + 2b \right].
%\end{align}
%CORRECT DERIVATIVE CHECK
%\begin{align}
%&\frac{1}{2}\left[ (x-b) \log |x+b| - \frac{x+b}{2} + 2b \right] + \frac{x+b}{2} \left[ \log |x+b| + \frac{x-b}{x+b} - \frac{1}{2} \right] = \\
%&= x\log |x+b| - \frac{x+b}{4} + b + \frac{x-b}{2} - \frac{x+b}{4} = x\log |x+b|.
%\end{align}
%%%%% FOUR INTEGRAL APPROACH
%We solve the four integrals separately
%\begin{align}
%I_1 &= \intx\int_{k_F - k}^{k_F} q \log |2k|\, \mathrm{d}q = \frac{k_F^2 - (k_F-k)^2}{2} \log |2k|,\\
%I_2 &= \intx\int_{k_F - k}^{k_F} q \log \left\vert q + \frac{k}{2} \right\vert \, \mathrm{d}q = \\
%&= \frac{q + k/2}{2} \klam\left[ (q - k/2) \log \left\vert q + \frac{k}{2} \right\vert - \frac{q + 3k/2}{2} \right]_{k_F - k}^{k_F} = \\
%&= \frac{q^2 - (k/2)^2}{2} \log \left\vert q + \frac{k}{2} \right\vert_{k_F - k}^{k_F} - \frac{(k_F + 3k/2)^2}{4} + \frac{(k_F + k/2)^2}{4} = \\
%&= \frac{k_F^2 - (k/2)^2}{2} \log \left\vert k_F + \frac{k}{2} \right\vert - \frac{(k_F - k)^2 - (k/2)^2}{2} \log \left\vert k_F - \frac{k}{2} \right\vert - \frac{k(k_F + k)}{2},\\
%I_3 &= \intx\int_{k_F - k}^{k_F} q \log |q - k_F| \, \mathrm{d}q = \frac{q - k_F}{2} \klam\left[ (q+k_F) \log |q-k_F| - \frac{q + 3k_F/2}{2} \right]_{k_F - k}^{k_F} = \\
%&= \frac{k}{2} \klam\left[ -k \log |k| - \frac{k_F/2 - k}{2} \right] = -\frac{k^2}{2}\log |k| + \frac{k^2}{4} - \frac{k_Fk}{8},\\
%I_4 &= \intx\int_{k_F - k}^{k_F} q \log |q + k_F| \, \mathrm{d}q = \frac{q + k_F}{2} \klam\left[ (q-k_F) \log |q+k_F| - \frac{q - 3k_F/2}{2} \right]_{k_F - k}^{k_F} = \\
%&= \frac{k_F^2}{4} - \frac{2k_F - k}{2} \klam\left[ -k\log |k| + \frac{k + k_F/2}{2} \right] = (k_F - k/2) k\log |k| + \frac{k^2}{4} - \frac{3k_Fk}{2}.
%\end{align}
%Before inserting these into our original expression note that
%\begin{align}
%I_3 + I_4 &= k(k_F-k) \log |k| + \frac{k^2}{2} - \frac{13k_Fk}{8},\\
%I_1 + I_2 &= \frac{2k_Fk - k^2}{2} \log |2k| + \frac{k_F^2 - (k/2)^2}{2} \log \left\vert \frac{k_F + k/2}{k_F - k/2} \right\vert - \frac{k^2 -k_Fk}{2} \log |k_F - k/2| - \frac{k(k_F + k)}{2}
%\end{align}
%and
%\begin{align}
%\sum I_j &= 2k(k_F - k)\log k -\frac{17k_Fk}{8} + \frac{2k_Fk-k^2}{2} \log 2 
%\end{align}
%%%% CORRECT SOLUTION USING ONLINE INTEGRATION
%\begin{align}
%&\frac{4x^2\ln\left(\frac{\left|k\right|\left|2x+k\right|}{\left|x-k_F\right|\left|x+k_F\right|}\right) - k^2\ln\left(\left|2x+k\right|\right)+2\left(2k_F^2\cdot\left(\ln\left(\left|x+k_F\right|\right)+\ln\left(\left|x-k_F\right|\right)\right)+x\cdot\left(x+k\right)\right)}{8}\bigg\vert_{k_F - k}^{k_F}\\
%&= \frac{4k_F^2 \log \frac{k(2k_F+k)}{2k_F} - k^2 \log (2k_F+k) + 4k_F^2 \log (2k_F) + k_F(k_F + k)}{8} - \\
%&\quad -\frac{4(k_F - k)^2 \log \frac{k(2k_F - k)}{k(2k_F - k)} - k^2 \log (2k_F - k) + 4k_F^2 \log (2k_F - k) + 4k_F^2 \log k + k_F (k_F - k)}{8}\\
%&= \frac{k_F^2}{2} \log |k(2k_F + k)| - \frac{k^2}{8} \log |2k_F + k| + k_F^2 + k_Fk + \\
%&\quad + \frac{k^2}{8} \log |2k_F - k| - \frac{k_F^2}{2} \log |2k_F - k| - \frac{k_F^2}{2} \log k - k_F^2 + k_Fk = \\
%&= \frac{k_F^2 - (k/2)^2}{2} \log \left\vert \frac{k_F + k/2}{k_F - k/2} \right\vert  + 2k_Fk.
%\end{align}
%%%% NEW ATTEMPT
%The first integral gives
%\begin{align}
%I_1 &= \int_{k_F-k}^{k_F} q\log |k(2q-k)| \, \mathrm{d}q = \intx\int_{k_F - k/2}^{k_F + k/2} (q + k/2) \log |2kq| \, \mathrm{d}q = \\
%&= \intx\int_{2k_Fk - k^2}^{2k_Fk + k^2} \frac{u}{(2k)^2} \log |u| \, \mathrm{d}u + k\frac{k}{2} \log |2k| + \intx\int_{k_F - k/2}^{k_F + k/2} \frac{k}{2} \log |q| \, \mathrm{d}q = \\
%&= \frac{2u^2\log |u| - u^2}{4(2k)^2} \Big\vert_{k(2k_F - k)}^{k(2k_F + k)} + \frac{k^2}{2} \log |2k| + \frac{k}{2} \klam\left[ q \log |q| - q \right]_{k_F - k/2}^{k_F + k/2} = \\
%&= \frac{(2k_F + k)^2 \kla\left( 2\log |k(2k_F + k)| - 1 \right) - (2k_F - k)^2 \kla\left( 2\log |k(2k_F - k)| - 1 \right)}{16} + \\
%&\quad + \frac{k^2}{2} \log |2k| + \frac{k}{2} \left[ k_F \log \left\vert \frac{k_F + k/2}{k_F - k/2} \right\vert + \frac{k}{2} \log |k_F^2 - (k/2)^2| - k \right] = \\
%&= -\frac{k_Fk}{4} + \frac{k_F^2 + (k/2)^2}{4} \log \left\vert \frac{k_F + k/2}{k_F - k/2} \right\vert + \frac{k_Fk}{4} \log \left\vert 4k^2 \left(k_F^2 - (k/2)^2\right) \right\vert + \dots = \\
%&= \frac{k_F^2 + (k/2)^2 + 2k_Fk}{4} \log\left\vert \frac{k_F + k/2}{k_F - k/2} \right\vert  + \frac{k(k_F + k)}{4} \left\vert 4k^2 \left(k_F^2 - (k/2)^2\right) \right\vert - \frac{k_Fk + 2k^2}{4}
%\end{align}
%The second integral gives
%\begin{align}
%I_2 &= \int_{k_F - k}^{k_F} q \log |q^2 - k_F^2|\, \mathrm{d}q = \intx\int_{0}^{k} (k_F - q) \log |q^2 - 2qk_F| \, \mathrm{d}q = \\
%&= \intx\int_{0}^{k} (k_F - q) \log |q| \, \mathrm{d}q + \intx\int_{0}^{k} (k_F - q) \log |q - 2k_F|\, \mathrm{d}q = \\
%&= k_Fk (\log k - 1) - \frac{q^2 (2\log |q| - 1)}{4} \Big\vert_0^k - \intx\int_{-2k_F}^{k-2k_F} q \log |q| \, \mathrm{d}q - k_F \intx\int_{0}^{k} \log |q - 2k_F|\, \mathrm{d}q
%\end{align}
%%%% NEW ATTEMPT
In this case we can actually find an anti-derivative
\begin{align}
I&= \int x\log \left\vert \frac{k(2x + k)}{x^2 - f^2} \right\vert = \\
&= \frac{x^2}{2} \log \left\vert \frac{k(2x + k)}{x^2 - f^2} \right\vert - \int \frac{x^2}{2} \frac{x^2 - f^2}{k(2x + k)} \klam\left[ \frac{2k}{x^2 - f^2} - \frac{k(2x + k)2x}{(x^2 - f^2)^2} \right] + C = \\
&= \frac{x^2}{2} \log \left\vert \frac{k(2x + k)}{x^2 - f^2} \right\vert  - \int \frac{x^2}{2x+k} + \int \frac{x^3}{x^2 - f^2} + C = \\
&= \frac{x^2}{2} \log \left\vert \frac{k(2x + k)}{x^2 - f^2} \right\vert - \frac{1}{8} \int \frac{(u - k)^2}{u} + \int \frac{z + f^2}{2z} + C = \\
&= \frac{x^2}{2} \log \left\vert \frac{k(2x + k)}{x^2 - f^2} \right\vert - \frac{1}{8} \klam\left[ \frac{u^2}{2} - 2ku + k^2 \log u \right] + \frac{z}{2} + \frac{f^2}{2} \log |z| + C.
\end{align}
The second integral was obtained by substituting $u := 2x + k$, going back to $x$ yields
\begin{align}
 - &\frac{\frac{(2x + k)^2}{2} - 2k(2x + k) + k^2 \log |2x + k|}{8} = \\
 &= \frac{-x^2}{4} + \frac{kx}{4} - \frac{k^2\log |2x + k|}{8} + C
\end{align}
Putting everything back together with $z := x^2 - f^2$ we find
\begin{align}
I &= \frac{x^2}{2} \log \left\vert \frac{k(2x + k)}{x^2 - f^2} \right\vert - \frac{ k^2 }{8}\log |2x + k|  + \frac{x(k+x)}{4} + \frac{f^2}{2} \log |x^2 - f^2| + C = \\
&= \frac{x^2}{2} \log |k| - \frac{x^2 - f^2}{2} \log |x^2 - f^2| + \frac{x^2 - (k/2)^2}{2}\log |2x+k| + \frac{x(k+x)}{4} + C.
\end{align}
When evaluating the integral at the bounds we may drop the arbitrary constant.
In our case $f=k_F$ and the bounds give
\begin{align*}
I\Big\vert_{k_F - k}^{k_F} &= \frac{k_F^2}{2}\log |k| + \frac{k_F^2 - (k/2)^2}{2} \log |2k_F + k| + \frac{k_F(k+k_F)}{4} - \\
&\quad -\frac{(k_F - k)^2}{2} \log |k| + \frac{k^2 - 2k_Fk}{2} \log |k^2 - 2k_Fk| - \\
&\quad - \frac{(k_F - k)^2 - (k/2)^2}{2} \log |2k_F - k| - \frac{k_F(k_F - k)}{4} = \numberthis \\
&= \frac{k_F^2 - (k/2)^2}{2} \log |2k_F + k| + \frac{k_Fk}{2} +\\
&\quad + \frac{k^2 - 2k_Fk - (k_F - k)^2 + (k/2)^2}{2} \log |2k_F - k| = \numberthis  \\
&= \frac{k_Fk}{2} + \frac{k_F^2 - (k/2)^2}{2} \log \left\vert \frac{k_F + k/2}{k_F - k/2} \right\vert. \numberthis 
\end{align*}
Inserting into the polarization we arrive at the Lindhard result
\begin{align}
\Pi(\textbf{k},0) &= \frac{4m/k}{(2\pi)^2} \klam\left[ \frac{k_Fk}{2} + \frac{k_F^2 - (k/2)^2}{2} \log \left\vert \frac{k_F + k/2}{k_F - k/2} \right\vert \right] = \\
&= \frac{2mk_F}{(2\pi)^2} \klam\left[ 1 + \frac{k_F^2 - (k/2)^2}{k_Fk} \log \left\vert \frac{k_F + k/2}{k_F - k/2} \right\vert \right].
\end{align}
%FIXME: Global minus sign missing, this \emph{has} to come from an error on the sheet (right?)\\
Note that we are missing a global minus sign in this result.\\
%Since we were also missing such a sign in the result for $\Pi(\textbf{k},\omega)$, maybe there is just a typo on the problem sheet in Equation (3), since we definitely need the minus sign in the Lindhard result to get the correct effective interaction potential.\\
The Dyson equation gives rise to an effective interaction $V^\tecxt\text{eff}(\textbf{k}, 0)$ of the form
\begin{align}
V^\tecxt\text{eff}(\textbf{k}, 0) &= \frac{V(\textbf{k}, 0)}{1 - V(\textbf{k}, 0) \Pi(\textbf{k}, 0)} = \frac{1}{k^2 + \frac{2mk_F}{(2\pi)^2} \left[ 1 + \frac{k_F^2 - (k/2)^2}{k_Fk} \log \left\vert \frac{k_F + k/2}{k_F - k/2} \right\vert \right]}.
\end{align}
In the long wavelength limit $k\rightarrow 0$ the polarization behaves like
\begin{align}
\Pi(\textbf{k},0) &= \frac{2mk_F}{(2\pi)^2} \klam\left[ 1 + \frac{k_F^2 - (k/2)^2}{k_Fk} \kla\left( 2\frac{k}{2k_F} + 2\frac{k^3}{3(2k_F)^3} + \mathcal{O}(k^5) \right) \right] = \\
&= \frac{2mk_F}{(2\pi)^2} \klam\left[ 1 + \kla\left( 1 - \frac{k^2}{4k_F^2} \right) \kla\left( 1 + \frac{k^2}{12k_F^2} + \mathcal{O}(k^4) \right) \right] = \\
&= \frac{2mk_F}{(2\pi)^2} \klam\left[ 2 - \frac{k^2}{6k_F^2} + \mathcal{O}(k^4) \right].
\end{align}
Inserting into the effective interaction yields
\begin{align}
V^\tecxt\text{eff}(\textbf{k},0) &= \frac{1}{k^2 \kla\left( 1 - \frac{m}{(2\pi)^2 3k_F} \right) + \frac{4mk_F}{(2\pi)^2} + \mathcal{O}(k^4)} \simeq \\
&\simeq \frac{1}{k^2 + \frac{4mk_F / (2\pi)^2}{1 - \frac{m}{(2\pi)^2 3k_F}}} =: \frac{1}{k^2 + k_\text{TF}^2}.
\end{align}
The Fourier transform of this potential is
\begin{align}
V^\tecxt\text{eff}(\textbf{r}) &= \intr\int_{\mathbb{R}^{3}} \frac{\e^{\i \textbf{k}\cdot \textbf{r}}}{k^2 + k_\text{TF}^2} \, \frac{\mathrm{d}^{3}k}{(2\pi)^3} = \\
&= \frac{1}{(2\pi)^2} \intx\int_{0}^{\infty} \intx\int_{-1}^{1} \frac{\e^{\i kr\cos\theta} }{k^2 + k_\text{TF}^2} \, \mathrm{d}(\cos\theta)\, k^2\mathrm{d}k = \\
&= \frac{1}{(2\pi)^2} \intx\int_{0}^{\infty} \frac{1}{\i kr} \frac{\e^{\i kr} - \e^{-\i kr}}{k^2 + k_\text{TF}^2} \, k^2\mathrm{d}k = \\
&= \frac{1}{(2\pi)^2 r} \tecxt\text{Im\,} \klammer\left\{ \intr\int_{\mathbb{R}^{ }} \frac{k\e^{\i kr}}{k^2 + k_\text{TF}^2} \, \mathrm{d}k \right\}.
\end{align}
The integrand has two poles at $k = \pm\i k_\text{TF}$, but only one of them lies above the real axis.
Since the integral along the arc vanishes, the residual theorem gives
\begin{align}
I &= 2\pi\i \lim _{k \rightarrow \i k_\tecxt\text{TF}} \frac{k\e^{\i kr}}{k + \i k_\text{TF}} = \pi\i \e^{-k_\text{TF}r}
\end{align}
and thus the effective interaction becomes
\begin{align}
V^\tecxt\text{eff}(\textbf{k},0) &= \frac{\e^{-k_\text{TF}r}}{4\pi r}.
\end{align}
The only issue is that the Thomas-Fermi wavevector was found to be
\begin{align}
k_\text{TF}^2 &= \frac{4mk_F}{(2\pi)^2} \frac{(2\pi)^2 3k_F}{(2\pi)^2 3k_F - m}
\end{align}
which is incompatible with the given $k_\text{TF}^2 = \frac{2mk_F}{(2\pi)^2}$ (for what it's worth, this is apparently the ``more precise'' result :) )

%\begin{thebibliography}{9}
%\bibitem{example} description, \url{https://www.exmapleurl}, day.\,month~2021.
%\end{thebibliography}

\end{document}